\section{The Propers: Detailed Commentary}
\subsection{A prayer from distress}
The entrance antiphon adapts Psalm 87:3 Vulgate, 88:3 Hebrew. Its speaker asks God to admit the prayer and attend to the plea. The surrounding psalm describes affliction, isolation and the nearness of death. Unlike a lament that ends by reporting deliverance, it leaves the sufferer's distress unresolved. The entrance therefore gives the assembly a petition to make; it does not announce that everyone praying has already recovered peace or health.

Augustine reads this psalm in the voice of Christ, whose human suffering is freely endured. His opening exposition distinguishes an apparent distance from God from a loss of divine grace. When the psalm asks that prayer come before God, Augustine understands hearing as merciful acceptance. This explains why the language of sorrow can belong to the Savior and to his members: the presence of grief does not itself prove that God has rejected the one praying.\footnote{Augustine, \textit{Expositions on the Psalms}, Psalm 87.1--4, headed Psalm LXXXVIII in NPNF1, vol. 8.}

The Collect asks for freedom directed toward God's service. Its concern extends to the person who serves, not merely to the removal of a particular inconvenience. Set beside the distressed entrance psalm, the petition can be understood as asking for the ability to respond to God amid the conditions of life. A promise that faithful people will encounter no adversity would contradict the lament with which they have entered.

\subsection{Wisdom sought and Wisdom seeking}
Wisdom 6 addresses rulers who must answer for their judgments. The surrounding exhortation makes the desire for wisdom a question of how power is exercised and how life is ordered. The appointed passage then describes Wisdom as radiant, accessible to love, and already approaching those who desire her. Its verbs concern both the seeker and Wisdom herself. A person rises early, thinks and watches; Wisdom shows herself, waits at the door and goes out to meet.

The historical English word ``preventeth'' in the study translation means that Wisdom comes before or anticipates the seeker. It does not mean that she obstructs the search. The final appointed verse develops that initiative further: Wisdom seeks persons and meets them in their ways. Its inclusion is especially important because the appointment's Wisdom 6:12--16 corresponds to 6:13--17 in the held Vulgate-based Challoner witness. Reading the unshifted numbers in that witness would omit Wisdom's own search.

Bonaventure explains the visibility of Wisdom through her own light. The seeker does not first become capable of seeing by an achievement independent of her. Love belongs to this perception, and seeking involves action as well as consideration. His exposition of the door describes an opening of intellect and will. The image therefore concerns the person who can understand and desire, not simply an intellectual puzzle solved from outside.\footnote{Bonaventure, \textit{Commentarius in librum Sapientiae}, chapter 6, Vulgate vv. 13--17, \textit{Opera omnia}, Quaracchi, vol. VI (1893), pp. 147--149.}

There is an instructive difference between Wisdom waiting at a door and the Gospel's wedding door being shut. The first passage presents the accessibility of the divine gift to a willing seeker. The parable describes the failure of preparation disclosed at an arrival. Together they invite an account of response that includes both God's prior generosity and a human opportunity that must be used. They do not require an identification of Wisdom's door with every detail of the wedding story.

\subsection{Thirst, praise and remembrance}
The Responsorial Psalm appoints Psalm 63:2--8 in Hebrew numbering, Psalm 62:2--8 in the Vulgate. The desire for God includes soul and flesh. Dryness is followed by recollection of divine power, praise, lifted hands and remembrance on the bed. The final appointed words speak of help and the shelter of God's wings. Danger remains in the whole psalm, but the later pursuit verses are outside this appointment.

Augustine extends the body's thirst toward its promised resurrection. The body belongs to the hope for God rather than being an obstacle that salvation must discard. He also distinguishes the soul's watchfulness from a refusal of bodily sleep. Remembering God at night does not condemn the ordinary rest that a body needs. His treatment of the raised hands makes another connection: the hands stretched toward God in prayer should also be stretched toward the needy in mercy.\footnote{Augustine, \textit{Expositions on the Psalms}, Psalm 62.1--12, especially the treatment of bodily thirst and raised hands; NPNF heading Psalm LXIII. The conjunction of prayer and mercy occurs in §10.}

Bellarmine explains the desire for God in dependence on help already received. His exposition is used here through John O'Sullivan's historical English abridgment, which omits some of the original's discussions. The practical consequence is consistent with the poem's sequence: the worshipper can remember God's support while still desiring what is not yet possessed. Neither thirst nor remembrance supplies a measure by which to rank another person's spiritual worth from outward emotional intensity.\footnote{Robert Bellarmine, \textit{A Commentary on the Book of Psalms}, O'Sullivan abridgment, Psalm 62, exposition of vv. 1--8.}

\subsection{The dead and the promised meeting}
Paul addresses a particular ignorance concerning those who have died. His reason for hope is Jesus' death and resurrection. The long form proceeds to the Lord's descent, the rising of the dead in Christ, the gathering of the living with them, and lasting life with the Lord. Its last direction makes this teaching something believers are to say to one another in consolation.

The short form ends after the foundation in Jesus' death and rising. It does not include the trumpet, the order of rising or the explicit command of mutual consolation at the end of the long form. When those details help explain the short form, they belong to the surrounding letter, not to words proclaimed within that shorter appointment. Both forms answer grief with Christ's resurrection; they differ in how much of Paul's explanation is heard.

Aquinas distinguishes sorrow for loss and separation from sorrow that has no hope. His explanation does not require a bereaved Christian to deny the pain of bodily absence. He argues that Christ's resurrection is both the cause and the pattern of the faithful person's resurrection. The promised meeting thus rests on something Christ has done, not on a mourner's ability to maintain an untroubled state of mind.\footnote{Thomas Aquinas, \textit{Super I ad Thessalonicenses}, chapter 4, lectio 2, on 1 Thessalonians 4:13--18.}

Chrysostom emphasizes that the dead lose no advantage to the living. He also distinguishes resurrection from the glory of the final meeting and explains the meeting through the welcome given to an arriving king. His treatment of mourning is markedly severe. Aquinas's explicit recognition of lawful grief should not be attributed backward to every phrase of that exhortation. Both expound hope in the resurrection, but their pastoral formulations of sorrow differ.\footnote{John Chrysostom, \textit{Homilies on First Thessalonians}, 6--8, opening expositions of 4:13, 4:13--15 and 4:15--18 respectively, NPNF1, vol. 13. Only these expository portions are used.}

\subsection{An unknown hour and a common sleep}
The acclamation selects and adapts Matthew 24:42a, 44. In the surrounding discourse, ignorance of the hour leads to a demand for continuing readiness. The entrusted servant must be found doing the work assigned. Chrysostom interprets this as service for the good of others: resources and authority should benefit those placed within one's care. The inability to calculate the arrival gives that service a continuing claim on the present.\footnote{John Chrysostom, \textit{Homilies on Matthew}, 77.1--3, NPNF1, vol. 10.}

The Gospel develops the difficulty through the delayed bridegroom. The foolish have taken lamps without the accompanying provision carried by the wise. At the cry, their lamps are going out; they seek oil, leave to obtain it, and return after the door is closed. The narrative gives an initial contrast in preparation, a common delay and sleep, and a final separation. Its warning is addressed to people who have already gone out to meet the bridegroom.

Augustine's explanation of sleep begins by testing a possible meaning against the whole story. If sleep meant charity growing cold, the wise would also have lost charity, since all sleep. Yet their entrance depends on perseverance. He therefore reads this common sleep as death and connects it with Paul's language about those who sleep. The concluding watchfulness belongs to faith, hope, charity and works while life continues.\footnote{Augustine, Sermon 93, §§6--9 and 17; NPNF Sermon 43.}

\subsection{The Passion and its reception}
The Prayer over the Offerings concerns the sacramental celebration of the Son's Passion and its inward reception. The relation between the action celebrated and the person participating matters to both interpretations that follow. In response to Christ's self-giving, a person may consider what it asks of conduct toward a neighbor, or examine whether participation is joined to inward allegiance.

\subsection{Two Communion alternatives}
The first Communion antiphon adapts the beginning of Psalm 22 Vulgate, Psalm 23 Hebrew. The speaker receives care from the shepherd: nourishment and refreshing water come from another's guidance. The whole psalm also speaks of the death-shadow, the table and the dwelling with God. Those later images help explain its confidence, but the appointed adaptation names the opening verses.

Augustine understands the shepherd as Christ, feeding the Church and leading it to baptismal refreshment. Bellarmine explains the promise of provision without making it a promise that the faithful will possess every worldly comfort. His point addresses a real difficulty in the poem: a confession of God's care can be made by someone who still experiences want or danger.\footnote{Augustine, \textit{Expositions on the Psalms}, Psalm 22.1--6, NPNF heading Psalm XXIII; Bellarmine, \textit{Commentary on the Psalms}, O'Sullivan abridgment, Psalm 22.}

The Luke alternative adapts the disciples' report of recognizing Jesus in the breaking of bread. Verse 35 belongs to their return and report in Jerusalem; the recognition it recounts occurred in the preceding encounter at Emmaus. Gregory follows that encounter through Christ's explanation of Scripture, the disciples' welcome to the apparent stranger, and their recognition at the meal. He argues that doing what they have heard enables their understanding. The hospitality is an act toward someone whose identity they have not yet grasped.\footnote{Gregory the Great, \textit{Homiliae in Evangelia}, 23.1--2, on Luke 24:13--35.}

These alternatives permit different emphases. The psalm speaks first of receiving a shepherd's care. Luke recalls an encounter in which disciples receive Christ as a guest and come to know him. Both can illuminate Communion, but the details of one should not disappear into a generalized account of nourishment.

\subsection{Thanksgiving and endurance}
The Prayer after Communion gives thanks and petitions for the Spirit's preservation of integrity. Reception is followed by a request concerning the continuing life of the recipient. This makes perseverance something for which the assembly asks divine help. The person who has received is still capable of needing correction, mercy and renewed purpose. The final petition can therefore accompany either the decision to offer practical help or the examination of the desire from which that help proceeds.
